\section{Introducción}

Una \textbf{red} es una estructura matemática compuesta por un conjunto de \textbf{vértices}
(nodos) y un conjunto de \textbf{aristas} (relaciones entre nodos). Formalmente, un grafo
$G = (V, E)$ donde $V$ es el conjunto de vértices y $E \subseteq V \times V$ el de aristas.

Las redes aparecen de forma natural en muchos dominios:

\begin{center}
\begin{tabular}{lll}
\toprule
\textbf{Dominio} & \textbf{Vértice} & \textbf{Arista} \\
\midrule
Biología & Proteína & Interacción física \\
Comercio & País & Flujo comercial \\
Cine & Actor & Aparición en la misma película \\
Informática & Dirección IP & Conexión de red \\
\bottomrule
\end{tabular}
\end{center}

Una pregunta central en el análisis de redes es: \textit{¿cuál elemento es el más
``importante''?} La respuesta depende de qué se entiende por importancia. Las
\textbf{métricas de centralidad} responden esta pregunta desde distintos ángulos,
cada una capturando una noción diferente de relevancia estructural.

En este proyecto se implementan \textbf{siete métricas de centralidad}, aplicadas
sobre dos redes reales: la red de interacciones proteína-proteína de levadura (\textit{Yeast PPI})
y la red de comercio internacional (\textit{Trade Network}).
